\section*{Chapter \ref{ch:rs:capacity-decay-and-crowding} --- Capacity, Decay and Crowding}

\begin{solution}{exo:rs:capacity-decay-and-crowding:1}
$k\sqrt{10^{10}} = 0.10$ gives $k = 10^{-6}$; $A^* = (2 \times 0.10/(3 \times 10^{-6}))^2 = \$4.44$ billion, where the net return is $0.10/3 = 3.3\%$.
\end{solution}

\begin{solution}{exo:rs:capacity-decay-and-crowding:2}
The return on capital falls from the first dollar, but the dollar profit is size times return: it keeps rising while the return falls more slowly than size
grows, up to the point where the marginal dollar earns nothing ($g - \frac32 k\sqrt A = 0$).
\end{solution}

\begin{solution}{exo:rs:capacity-decay-and-crowding:3}
A position of 2\% of \$2 billion, \$40 million, sold over three days: \$13.3 million a day against \$200 million traded, 6.7\% of volume. The push is
$0.7 \times 2\% \times \sqrt{0.067} = 0.36\%$ a day.
\end{solution}

\begin{solution}{exo:rs:capacity-decay-and-crowding:4}
Each name's push is proportional to the square root of the seller's trade in it, and the trade is proportional to the fraction sold: the loss of a fund holding
the same name scales as its square root. A non-seller's loss is its book times the push, so it is proportional to how much of its book lies where the seller
trades: the overlap.
\end{solution}

\begin{solution}{exo:rs:capacity-decay-and-crowding:5}
The loss was confined to long--short factor books, not the market: someone was selling the long side and buying back the short side of value and momentum books
at the same time. That is the signature of deleveraging by funds holding similar books, as Khandani and Lo concluded, not of news.
\end{solution}

\begin{solution}{exo:rs:capacity-decay-and-crowding:6}
Costs grow faster than size, so the damage of size comes from trading. A book re-optimised at each size trades less as it grows (at \$10 billion it turns over 1.8\% of
capital a day against 227\% at \$10 million in chapter~27), which keeps its costs from outrunning its return; the signal was the same.
\end{solution}

\begin{solution}{exo:rs:capacity-decay-and-crowding:7}
With identical books, day fifteen over the peak is $1.30/3.07 = 0.42$ for the whole book and $0.41/0.97 = 0.42$ for a tenth: above the permanent share of 0.3
because the peak, on the third day, already includes some decay of the first days' \omterm{def:rs:transaction-costs-in-research:impact}{temporary impact}.
\end{solution}

\begin{solution}{exo:rs:capacity-decay-and-crowding:8}
A positive Sharpe ratio is not a \omterm{def:rs:capacity-decay-and-crowding:capacity}{capacity}: at \$5 billion the fixed rule's ratio has fallen to a fifth of its best (0.34 against 1.79) and its dollar profit is
past its peak (\$146 million at \$3 billion, \$122 million at \$5 billion). By the chapter's criterion its \omterm{def:rs:capacity-decay-and-crowding:capacity}{capacity} is \$2.02 billion, and the estimate is itself
uncertain; run it below that.
\end{solution}

\begin{solution}{pb:rs:capacity-decay-and-crowding:1}
\begin{enumerate}
\item $A^* = (2g/3k)^2$, where $g - \frac32 k\sqrt A = 0$; the net return there is $g/3$.
\item 1.79, 1.74, 1.64, 1.48, 1.18, 0.90, 0.68, 0.34, $-0.28$, $-1.85$ at \$10m, \$30m, \$100m, \$300m, \$1bn, \$2bn, \$3bn, \$5bn, \$10bn, \$30bn.
\item Profit peaks at \$3 billion (\$146 million a year); the Sharpe ratio halves at \$2.02 billion.
\item The book trades less as it grows: Sharpe ratio 2.43 at \$1 billion, 1.36 at \$30 billion, with profit still growing (\$806 million).
\item Below the \omterm{def:rs:capacity-decay-and-crowding:capacity}{profit-maximising size}, where the curve is flat, because the curve will move (decay, competitors) and the estimate is uncertain.
\item The signal decays as others find it, and the costs rise as others trade the same names.
\item \omterm{def:rs:capacity-decay-and-crowding:crowding}{Crowding} (an arbitrageur cannot know how many peers are in the same trade) and leverage (a private choice that creates a fire-sale externality).
\item The average abnormal return correlation among a momentum strategy's stocks; high values predicted strong reversals of momentum returns.
\item Correlation among the strategy's own names net of factors, overlap with estimated books of others, short interest and days to cover, and the book's liquidity.
\item Five funds of \$2 billion, gross four times capital, 100 names a side from 500 stocks trading \$200 million a day at 2\% volatility; one sells over three days;
square-root impact with $\eta = 0.7$, 30\% permanent, the rest decaying with a half-life of one day.
\item From $-0.03\%$ (10\% sold, overlap 0.04) to $-3.07\%$ (all sold, identical books); the table in the text.
\item On the third day, the last of the selling; at day fifteen the whole-book case has recovered to $-1.30\%$.
\item From 6 to 9 August value fell 2.28\% and momentum 3.43\% (10.6 and 12.5 daily standard deviations) while the market rose 1.39\%; on 10 August value regained
1.42\% and momentum 1.15\%.
\item Losses of long--short value-type books that began in July, two \omterm{def:rs:capacity-decay-and-crowding:crowding}{unwinds} on 1 and 6 August starting with financials, long book-to-market and short momentum,
and a withdrawal of market-making capital from 8 August.
\item \textbf{Named result.} With identical books, the non-selling funds' peak loss is 0.97\%, 1.54\%, 2.17\% and 3.07\% of capital when a tenth, a quarter, half and
all of the book is sold (proportionally less at lower overlap); the fixed rule's net Sharpe ratio halves at \$2.02 billion.
\item Not sell at the bottom if it can hold: the temporary part comes back; reduce only as much as its own limits require, and provide liquidity if it has capital.
\item They turn losses into more selling: the peaks deepen and spread to funds that did not overlap with the first seller.
\item Its strategies' trades in the same names add up under a concave cost law, so their combined \omterm{def:rs:capacity-decay-and-crowding:capacity}{capacity} is less than the sum.
\item The \omterm{def:rs:capacity-decay-and-crowding:capacity}{capacity curve} at the firm's costs, the turnover and holding period, the names' liquidity, and who else is likely to hold the same book.
\item Its costs as it grows, set by how much it trades in which names, and the other capital in the same trade.
\end{enumerate}
\end{solution}

\begin{solution}{iq:rs:capacity-decay-and-crowding:1}
In the week of 6 August 2007 long--short quantitative equity books lost heavily with no news, as value and momentum factors fell by ten or more daily standard
deviations while the market rose; on 10 August much of it came back. Khandani and Lo attribute it to the unwinding of similar books by several funds and a
withdrawal of market-making capital.
\end{solution}

\begin{solution}{iq:rs:capacity-decay-and-crowding:2}
Re-run the \omterm{def:rs:vectorised-backtests:backtest}{backtest} at a grid of sizes with the firm's fitted cost model (spread and square-root impact), re-optimised at each size; read the net Sharpe ratio and
dollar profit against size; take the \omterm{def:rs:capacity-decay-and-crowding:capacity}{profit-maximising size} and a criterion such as the size at which the net Sharpe ratio halves; redraw with live data.
\end{solution}

\begin{solution}{iq:rs:capacity-decay-and-crowding:3}
Look for excess correlation among the strategy's names net of risk-model factors (\omterm{def:rs:capacity-decay-and-crowding:crowding}{comomentum}), overlap with estimated books of others (holdings filings, returns
of similar products), high short interest and days to cover on the short side, and whether the strategy's losses coincide with others'.
\end{solution}

\begin{solution}{iq:rs:capacity-decay-and-crowding:4}
Estimate the overlap and the seller's size; expect a temporary loss on the overlapping part; do not sell into it unless limits force you; reduce the most
illiquid overlapping positions early if you must; consider providing liquidity to the seller if the permanent part is small.
\end{solution}

\begin{solution}{iq:rs:capacity-decay-and-crowding:5}
Profit is size times net return: the return (and the Sharpe ratio) falls from the first dollar under concave costs, but profit rises while size grows faster than
the return falls, up to $A^* = (2g/3k)^2$.
\end{solution}

\begin{solution}{iq:rs:capacity-decay-and-crowding:6}
If the seller trades $fQ_i$ of each name over a day and the push is $\eta\sigma_i\sqrt{fQ_i/V_i}$, a fund holding $h_i$ of the name loses
$\sum_i h_i\eta\sigma_i\sqrt{fQ_i/V_i} \propto \sqrt f$ on the day, of which the permanent share stays; spreading the sale over $d$ days and letting the rest decay
gives the chapter's paths.
\end{solution}
